\subsection{环的幂零元的刻画}

命题1. --- 设 A 是一个交换环，x 是 A 的一个元素。为使 x 是幂零的，充要条件是 1 - xT 在环 A{[}T{]} 中可逆。

我们注意到 A{[}T{]} 是形式幂级数环 A{[}{[}T{]}{]} 的一个子环，且 1 - xT 在 A{[}{[}T{]}{]} 中有逆元 \(\sum_{n=0}^{\infty} x^{n} T^{n}\) （IV，第30页，命题5）。为使 1 - xT

在 A{[}T{]} 中可逆，充要条件是 \(\sum_{n=0}^{\infty} x^n T^n\) 是一个多项式，

即 \(x\) 是幂零的。

命题2. --- 设 A 是一个交换环。 A 的幂零元集合是 A 的一个理想，等于 A 的所有素理想的交集。

设 \(x\) 是 \(\mathbf{A}\) 的一个幂零元， \(\mathfrak{p}\) 是 \(\mathbf{A}\) 的一个素理想。 \(x \mod \mathfrak{p}\) 的剩余类是整环 \(\mathbf{A} / \mathfrak{p}\) 的一个幂零元，因而是零；于是我们有 \(x \in \mathfrak{p}\) 。

设 \(x\) 是 \(A\) 的一个非幂零元。由命题1，主理想 \((1 - xT)\) （ \(A[T]\) 的）异于 \(A[T]\) 。由 Krull 定理（I，第104页），存在一个极大理想 \(m\) （ \(A[T]\) 的），它包含 \(1 - xT\) 。于是 \(m\) 是 \(A[T]\) 的素理想，因此 \(p = A \cap m\) 是 \(A\) 的素理想。我们有 \(1 \notin m\) 且 \(1 - Tx \in m\) ，因此 \(Tx \notin m\) ，从而 \(x \notin p\) 。

于是我们证明了，A 的幂零元集合 n 是 A 的所有素理想的交集；由于理想的任意交仍是理想，故 n 是一个理想。

推论. --- 为使一个交换环是既约的（V，第34页，定义2），充要条件是它同构于一个域的乘积的子环。

该条件显然是充分的。

设 A 是既约的。由命题2，A 的素理想集合的交 n 化为 0。对 A 的每个素理想 p，设 \(k(\mathfrak{p})\) 是 A/p 的分式域， \(\varphi_{p}\) 是 A 到 \(k(\mathfrak{p})\) 的典范同态。设 \(\varphi\) 是 A 到 \(\prod_{\mathfrak{p}} k(\mathfrak{p})\) 的同态，其指标为 p 的分量是 \(\varphi_{p}\) 。 \(\varphi_{\mathfrak{p}}\) 的核是 \(\mathfrak{p}\) ，故 \(\varphi\) 的核是 \(\mathfrak{n} = 0\) ；因此 \(\varphi\) 是 \(\mathbf{A}\) 到 \(\prod_{\mathfrak{p}} k(\mathfrak{p})\) 的一个子环上的同构。

